% =============================================================================
% paper62_automorphe_formen_de.tex
% Paper 62 (DE): Automorphe Formen — Brücke zwischen Zahlentheorie und Darstellungstheorie
% Autor: Michael Fuhrmann
% Projekt: Specialist Mathematics Project, Build 143
% Datum: 2026-03-12
% =============================================================================

\documentclass[12pt]{amsart}

\usepackage{amsmath, amssymb, amsthm}
\usepackage{mathtools}
\usepackage{hyperref}
\usepackage{geometry}
\usepackage{enumitem}
\usepackage[T1]{fontenc}
\usepackage[utf8]{inputenc}
\usepackage[ngerman]{babel}

\geometry{margin=2.5cm}

% Theorem-Umgebungen (Deutsch)
\newtheorem{theorem}{Satz}[section]
\newtheorem{lemma}[theorem]{Lemma}
\newtheorem{proposition}[theorem]{Proposition}
\newtheorem{korollar}[theorem]{Korollar}
\newtheorem{conjecture}{Vermutung}[section]
\theoremstyle{definition}
\newtheorem{definition}[theorem]{Definition}
\newtheorem{beispiel}[theorem]{Beispiel}
\newtheorem{bemerkung}[theorem]{Bemerkung}

% Nützliche Makros
\newcommand{\HH}{\mathbb{H}}
\newcommand{\ZZ}{\mathbb{Z}}
\newcommand{\QQ}{\mathbb{Q}}
\newcommand{\RR}{\mathbb{R}}
\newcommand{\CC}{\mathbb{C}}
\newcommand{\NN}{\mathbb{N}}
\newcommand{\AA}{\mathbb{A}}
\newcommand{\FF}{\mathbb{F}}
\newcommand{\GL}{\mathrm{GL}}
\newcommand{\SL}{\mathrm{SL}}
\newcommand{\PSL}{\mathrm{PSL}}
\newcommand{\Gal}{\mathrm{Gal}}
\newcommand{\vol}{\mathrm{vol}}
\newcommand{\tr}{\mathrm{tr}}
\newcommand{\im}{\mathrm{Im}}
\newcommand{\re}{\mathrm{Re}}
\newcommand{\abs}[1]{\left|#1\right|}
\newcommand{\norm}[1]{\left\|#1\right\|}
\newcommand{\inner}[2]{\left\langle #1,\, #2 \right\rangle}
\newcommand{\floor}[1]{\left\lfloor #1 \right\rfloor}

% -----------------------------------------------------------------------
\title{Automorphe Formen:\\
Brücke zwischen Zahlentheorie und Darstellungstheorie}

\author{Michael Fuhrmann}

\address{Specialist Mathematics Project, Build 143}

\email{specialist-maths@project.local}

\date{2026-03-12}

\keywords{Automorphe Formen, Maass-Formen, Spektraltheorie, Eisenstein-Reihen,
Langlands-Programm, Selberg-Spurformel, Hecke-Operatoren, kuspidales Spektrum}

\subjclass[2020]{11F70, 11F72, 11F37, 11F41, 22E55}

% -----------------------------------------------------------------------
% Abstract VOR \maketitle (amsart-Anforderung)
\begin{abstract}
Automorphe Formen gehören zu den tiefgründigsten und weitreichendsten Strukturen
der modernen Mathematik: Sie verbinden Zahlentheorie, harmonische Analysis und
Darstellungstheorie auf einzigartige Weise. In dieser Arbeit entwickeln wir die
Theorie systematisch — von den klassischen Modulformen auf $\SL(2,\ZZ)$ bis zur
adèlischen Formulierung auf $\GL(2,\AA)$. Wir behandeln Maass-Formen und die
Spektralzerlegung von $L^2(\Gamma \backslash \HH)$, beweisen die Selberg-Spurformel
(bewiesen) und diskutieren die Selberg-$\tfrac{1}{4}$-Vermutung (offen). Wir
analysieren Eisenstein-Reihen und die Langlands-Spektralzerlegung
$L^2(\Gamma \backslash \HH) = L^2_{\mathrm{cusp}} \oplus L^2_{\mathrm{kont}}$
(bewiesen von Langlands, 1976). Als Anwendungen stellen wir die Voronoi-Summenformel
(bewiesen), Subkonvexitätsschranken und den Sato-Tate-Satz (bewiesen von Taylor u.\,a.,
2008) vor. Abschließend werden neuere Entwicklungen zusammengefasst: die Langlands-Funk\-torialität
für $\GL(2) \to \GL(n)$ (teilweise bewiesen) und Ngô Bảo Châus Beweis des Fundamentallemmas
(Fields-Medaille 2010, bewiesen).
\end{abstract}

\maketitle

% =============================================================================
\section{Einleitung}
\label{sec:einleitung}
% =============================================================================

Die Theorie der automorphen Formen entstand aus der klassischen Untersuchung von
Modulformen im neunzehnten Jahrhundert und erhielt ihre moderne Gestalt durch die
bahnbrechende Arbeit von Langlands in den 1960er und 1970er Jahren. Heute sind
automorphe Formen die zentralen Objekte eines der ehrgeizigsten Programme der
Mathematik: des Langlands-Programms, das Zahlentheorie, algebraische Geometrie
und Darstellungstheorie durch ein Netz tiefer Korrespondenzen zu vereinen versucht.

Auf der klassischen Ebene ist eine \emph{Modulform} vom Gewicht $k$ und Niveau $N$
eine holomorphe Funktion $f: \HH \to \CC$ auf der oberen Halbebene, die das
Transformationsgesetz
\[
f\!\left(\frac{az+b}{cz+d}\right) = (cz+d)^k\, f(z)
\quad \text{für alle } \begin{pmatrix}a&b\\c&d\end{pmatrix} \in \Gamma_0(N)
\]
erfüllt und geeignete Wachstumsbedingungen an den Spitzen besitzt. Das Paradebeis\-piel
ist Ramanujans Deltafunktion
\[
\Delta(z) = q \prod_{n=1}^\infty (1-q^n)^{24} = \sum_{n=1}^\infty \tau(n)\, q^n,
\quad q = e^{2\pi i z}.
\]
Sie ist die einzige (bis auf Skalierung) normierte kuspidierte Neuform vom Gewicht 12
und Niveau 1.

Die moderne Sichtweise, die auf Langlands, Jacquet, Godement-Jacquet zurückgeht,
hebt diese Objekte zu \emph{automorphen Darstellungen}: irreduziblen Darstellungen
von $\GL(2,\AA)$ (wobei $\AA = \AA_\QQ$ der Adèlring ist), die im Raum der
automorphen Formen auf $\GL(2,\QQ) \backslash \GL(2,\AA)$ realisiert werden.

\medskip
\noindent\textbf{Gliederung.}
Abschnitt~\ref{sec:klassisch} bespricht klassische Modulformen.
Abschnitt~\ref{sec:GL2} entwickelt die adèlische Formulierung auf $\GL(2)$.
Abschnitt~\ref{sec:maass} behandelt Maass-Formen.
Abschnitt~\ref{sec:spektral} entwickelt die Spektraltheorie.
Abschnitt~\ref{sec:langlands} gibt eine Vorschau auf das Langlands-Programm.
Abschnitt~\ref{sec:eisenstein} untersucht Eisenstein-Reihen.
Abschnitt~\ref{sec:kuspidal_eis} bespricht kuspidales versus Eisenstein-Spektrum.
Abschnitt~\ref{sec:anwendungen} stellt Anwendungen vor.
Abschnitt~\ref{sec:neueres} fasst neuere Entwicklungen zusammen.

% =============================================================================
\section{Klassische Modulformen: Eine Übersicht}
\label{sec:klassisch}
% =============================================================================

\subsection{Die obere Halbebene und \texorpdfstring{$\SL(2,\ZZ)$}{SL(2,Z)}}

Sei $\HH = \{z \in \CC : \im(z) > 0\}$ die \emph{Poincarésche obere Halbebene},
ausgestattet mit der hyperbolischen Metrik $ds^2 = y^{-2}(dx^2 + dy^2)$ für
$z = x+iy$. Die Gruppe $\SL(2,\RR)$ wirkt auf $\HH$ durch Möbius-Transformationen:
\[
\gamma \cdot z = \frac{az + b}{cz + d}, \quad \gamma = \begin{pmatrix}a&b\\c&d\end{pmatrix}
\in \SL(2,\RR).
\]
Der Kern dieser Wirkung ist $\{\pm I\}$; die effektiv wirkende Gruppe ist
$\PSL(2,\RR) = \SL(2,\RR)/\{\pm I\}$.

Die Modulgruppe $\Gamma = \SL(2,\ZZ)$ wirkt eigentlich diskontinuierlich auf $\HH$,
und der Quotient $\Gamma \backslash \HH$ ist eine nicht-kompakte Riemannsche Fläche
vom Geschlecht null mit einer Spitze (bei $i\infty$). Ein Fundamentalbereich ist
\[
\mathcal{F} = \left\{ z \in \HH : \abs{z} \geq 1,\; \abs{\re(z)} \leq \tfrac{1}{2} \right\}.
\]

\subsection{Modulformen und der Strich-Operator}

\begin{definition}
Für $\gamma = \bigl(\begin{smallmatrix}a&b\\c&d\end{smallmatrix}\bigr) \in \GL^+(2,\RR)$
und $k \in \ZZ$ ist der \emph{Strich-Operator vom Gewicht $k$}:
\[
(f \mid_k \gamma)(z) = \det(\gamma)^{k/2}\, (cz+d)^{-k}\, f(\gamma z).
\]
Eine \emph{Modulform vom Gewicht $k$ und Niveau $N$} ist eine holomorphe Funktion
$f : \HH \to \CC$, sodass:
\begin{enumerate}[label=(\roman*)]
\item $f \mid_k \gamma = f$ für alle $\gamma \in \Gamma_0(N)$;
\item $f$ holomorph in jeder Spitze von $\Gamma_0(N)$ ist.
\end{enumerate}
Verschwindet $f$ zusätzlich in jeder Spitze, so heißt $f$ \emph{Spitzenform} (Kuspform).
\end{definition}

Der Raum der Modulformen vom Gewicht $k$ und Niveau $N$ wird mit $M_k(\Gamma_0(N))$
bezeichnet, der Unterraum der Spitzenformen mit $S_k(\Gamma_0(N))$. Beide sind
endlichdimensionale $\CC$-Vektorräume. Das Riemann-Roch-Theorem liefert die Dimensionsformel:
\[
\dim S_k(\Gamma_0(N)) = (k-1)(g-1) + \floor{k/4}\,\nu_2 + \floor{k/3}\,\nu_3 + \tfrac{k}{2}\,\nu_\infty
\quad (k \geq 2 \text{ gerade}),
\]
wobei $g$ das Geschlecht ist und $\nu_2, \nu_3, \nu_\infty$ elliptische Punkte
und Spitzen zählen.

\subsection{Hecke-Operatoren}

Für eine Primzahl $p \nmid N$ wirkt der \emph{Hecke-Operator} $T_p$ auf
$M_k(\Gamma_0(N))$ durch
\[
(T_p f)(z) = p^{k-1}\, f(pz) + \frac{1}{p} \sum_{j=0}^{p-1} f\!\left(\frac{z+j}{p}\right).
\]
Die Hecke-Operatoren $\{T_n\}_{(n,N)=1}$ bilden eine kommutative Algebra — die
\emph{Hecke-Algebra} — und $S_k(\Gamma_0(N))$ besitzt eine Orthonormalbasis
simultaner Eigenfunktionen (\emph{Hecke-Eigenformen}), die \emph{Neuformen} heißen.
Ist $f = \sum a_n q^n$ eine normierte Neuform mit $a_1 = 1$, so gilt
\[
T_p f = a_p f \quad \text{und} \quad L(f,s) = \sum_{n=1}^\infty \frac{a_n}{n^s}
= \prod_{p \nmid N} \frac{1}{1 - a_p p^{-s} + p^{k-1-2s}}
\cdot \prod_{p \mid N} \frac{1}{1 - a_p p^{-s}}.
\]

% =============================================================================
\section{Automorphe Formen auf \texorpdfstring{$\GL(2)$}{GL(2)}}
\label{sec:GL2}
% =============================================================================

\subsection{Der Adèlring}

Der \emph{Adèlring} $\AA = \AA_\QQ$ ist das eingeschränkte direkte Produkt
\[
\AA = \RR \times \prod_{p}' \QQ_p = \left\{ (x_\infty, x_2, x_3, x_5, \ldots) : x_p \in \ZZ_p \text{ für fast alle } p \right\},
\]
ausgestattet mit der eingeschränkten Produkttopologie. Die diagonale Einbettung
$\QQ \hookrightarrow \AA$ ist diskret, und $\QQ \backslash \AA$ ist kompakt.

Die Gruppe $\GL(2,\AA)$ enthält $\GL(2,\QQ)$ als diskrete Untergruppe (via der
diagonalen Einbettung), und der Doppelnebenklassenraum
$\GL(2,\QQ) \backslash \GL(2,\AA) / \GL(2,\hat\ZZ)$ wird natürlich mit einer
disjunkten Vereinigung von Modulkurven $\Gamma_0(N) \backslash \HH^*$ identifiziert.

\subsection{Automorphe Darstellungen}

\begin{definition}
Eine \emph{automorphe Form auf $\GL(2)$} ist eine glatte Funktion
$\phi: \GL(2,\AA) \to \CC$ mit folgenden Eigenschaften:
\begin{enumerate}[label=(\roman*)]
\item \textbf{Linksinvarianz}: $\phi(\gamma g) = \phi(g)$ für alle $\gamma \in \GL(2,\QQ)$;
\item \textbf{Zentral-Charakter}: $\phi(zg) = \omega(z)\phi(g)$ für alle $z \in \AA^\times$,
  wobei $\omega: \QQ^\times \backslash \AA^\times \to \CC^\times$ ein Hecke-Charakter ist;
\item \textbf{$K$-Endlichkeit}: die Bahn unter der maximalen kompakten Untergruppe
  $K = \mathrm{O}(2) \times \prod_p \GL(2,\ZZ_p)$ erzeugt einen endlichdimensionalen Raum;
\item \textbf{$\mathfrak{z}$-Endlichkeit}: $\phi$ wird durch ein Ideal endlicher Kodimension
  im Zentrum der universellen einhüllenden Algebra annulliert;
\item \textbf{Moderates Wachstum}: es gibt $C, A > 0$ mit $|\phi(g)| \leq C \norm{g}^A$.
\end{enumerate}
\end{definition}

Eine \emph{automorphe Darstellung} $\pi$ ist ein irreduzibler Bestandteil der
rechtsseitigen regulären Darstellung von $\GL(2,\AA)$ auf dem Raum der automorphen Formen.

\subsection{Lokale-Globale-Zerlegung}

\begin{theorem}[Flath, 1979 — bewiesen]
Jede irreduzible automorphe Darstellung $\pi$ von $\GL(2,\AA)$ zerlegt sich als
eingeschränktes Tensorprodukt
\[
\pi \cong \pi_\infty \otimes \bigotimes_p' \pi_p,
\]
wobei $\pi_\infty$ ein zulässiges $(\mathfrak{g}_\infty, K_\infty)$-Modul ist und
$\pi_p$ eine irreduzible zulässige Darstellung von $\GL(2,\QQ_p)$ ist, die für
fast alle $p$ unverzweigt ist (d.\,h.\ $\pi_p^{\GL(2,\ZZ_p)} \neq 0$).
\end{theorem}

Die archimedische Komponente $\pi_\infty$ bestimmt das Gewicht (für holomorphe Formen)
oder den Eigenwert des Laplace-Operators (für Maass-Formen). Die nicht-archimedischen
Komponenten $\pi_p$ kodieren die Hecke-Eigenwerte $a_p$.

% =============================================================================
\section{Maass-Formen}
\label{sec:maass}
% =============================================================================

\subsection{Der hyperbolische Laplace-Operator}

Der \emph{hyperbolische Laplace-Beltrami-Operator} auf $\HH$ ist
\[
\Delta = -y^2 \!\left(\frac{\partial^2}{\partial x^2} + \frac{\partial^2}{\partial y^2}\right).
\]
Dieser Operator ist selbstadjungiert auf $L^2(\Gamma \backslash \HH, d\mu)$
bezüglich des invarianten Maßes $d\mu = y^{-2} dx\, dy$ und vertauscht mit der
$\SL(2,\RR)$-Wirkung. Sein Spektrum auf $\Gamma \backslash \HH$ besteht aus
$\{0\}$ (Konstanten), einem diskreten Teil $\{\lambda_n\}$ (mit
$0 < \lambda_1 \leq \lambda_2 \leq \cdots$) und einem stetigen Teil $[1/4, \infty)$.

\subsection{Maass-Wellenformen}

\begin{definition}
Eine \emph{Maass-Form} vom Gewicht $0$ und Niveau $N$ ist eine glatte Funktion
$u : \HH \to \CC$ mit:
\begin{enumerate}[label=(\roman*)]
\item $u(\gamma z) = u(z)$ für alle $\gamma \in \Gamma_0(N)$;
\item $\Delta u = \lambda u$ für einen Eigenwert $\lambda \in \RR$;
\item $u \in L^2(\Gamma_0(N) \backslash \HH)$ (oder höchstens polynomiales Wachstum).
\end{enumerate}
Verschwindet der konstante Fourier-Koeffizient $\int_0^1 u(x+iy)\, dx = 0$ für alle
$y > 0$, so heißt $u$ eine \emph{Maass-Spitzenform}.
\end{definition}

\begin{bemerkung}
Der Eigenwert $\lambda$ erfüllt $\lambda \geq 0$ (wegen Selbstadjungiertheit) und
lässt sich als $\lambda = \frac{1}{4} + r^2$ mit $r \in \RR \cup i[-\tfrac{1}{2}, \tfrac{1}{2}]$
schreiben. Der Fall $r \in \RR$ (d.\,h.\ $\lambda \geq 1/4$) entspricht temperierten Darstellungen.
\end{bemerkung}

\subsection{Fourier-Entwicklung von Maass-Formen}

Eine Maass-Spitzenform $u$ mit Eigenwert $\lambda = \frac{1}{4}+r^2$ besitzt die
Fourier-Entwicklung
\[
u(x+iy) = \sum_{n \neq 0} a_n\, \sqrt{y}\, K_{ir}(2\pi|n|y)\, e^{2\pi i n x},
\]
wobei $K_{ir}$ die modifizierte Besselfunktion zweiter Art (Macdonald-Funktion) ist.
Das klassische Beispiel ist die \emph{nichts-holomorphe Eisenstein-Reihe}
$E(z,s) = \sum_{\gamma \in \Gamma_\infty \backslash \Gamma} \im(\gamma z)^s$,
die eine Eigenfunktion von $\Delta$ mit Eigenwert $s(1-s)$ ist.

\subsection{Hecke-Theorie für Maass-Formen}

Die Hecke-Operatoren $T_n$ (für $(n,N)=1$) wirken auf Maass-Formen genauso wie
auf holomorphe Modulformen. Für eine simultane Hecke-Eigenfunktion $u$ mit
$T_n u = \lambda_n u$ ist die zugehörige $L$-Funktion
\[
L(u, s) = \sum_{n=1}^\infty \frac{a_n}{n^s} = \prod_p \frac{1}{1 - a_p p^{-s} + p^{-2s}}
\]
analytisch fortsetzbar und erfüllt die Funktionalgleichung $s \leftrightarrow 1-s$.

% =============================================================================
\section{Spektraltheorie}
\label{sec:spektral}
% =============================================================================

\subsection{Spektralzerlegung von \texorpdfstring{$L^2(\Gamma \backslash \HH)$}{L2}}

Das Spektrum von $\Delta$ auf $L^2(\Gamma \backslash \HH)$ besteht aus zwei Teilen:
\begin{itemize}
\item \textbf{Diskretes Spektrum}: Eigenwerte $0 = \lambda_0 < \lambda_1 \leq \lambda_2 \leq \cdots$
  entsprechend der konstanten Funktion (für $\lambda_0 = 0$) und Maass-Spitzenformen
  (für $\lambda_j > 0$); diese besitzen $L^2$-Eigenfunktionen.
\item \textbf{Stetiges Spektrum}: das Intervall $[1/4, \infty)$, parametrisiert durch
  Eisenstein-Reihen $E(z, \tfrac{1}{2}+it)$ für $t \in \RR$.
\end{itemize}

\subsection{Die Selberg-Spurformel}

\begin{theorem}[Selberg, 1956 — bewiesen]
\label{satz:selberg_spur}
Sei $h$ eine gerade holomorphe Funktion im Streifen $|\im(r)| < \frac{1}{2}+\epsilon$
mit $h(r) = O((1+|r|)^{-2-\epsilon})$, und sei $g(u) = \frac{1}{2\pi} \int_\RR h(r) e^{iru}\, dr$.
Dann gilt:
\begin{align*}
\sum_j h(r_j) + \frac{1}{4\pi}\int_{-\infty}^\infty h(r) \frac{-\varphi'}{\varphi}\!\left(\tfrac{1}{2}+ir\right) dr
&= \frac{\vol(\Gamma \backslash \HH)}{4\pi}\int_\RR h(r)\, r \tanh(\pi r)\, dr \\
&+ \sum_{\{P\}_{\Gamma,\text{hyp}}} \frac{\log N(P_0)}{N(P)^{1/2} - N(P)^{-1/2}}\, g(\log N(P)) \\
&+ \text{(elliptische Terme)} + \text{(parabolische Terme)},
\end{align*}
wobei die linke Summe über alle $r_j$ mit $\lambda_j = \frac{1}{4}+r_j^2$ läuft,
$\varphi(s)$ die Streumatrix ist, und die rechte Seite über $\Gamma$-Konjugationsklassen
(Identität, hyperbolisch, elliptisch, parabolisch) summiert.
\end{theorem}

Die Selberg-Spurformel ist eine tiefgründige Identität, die die Spektralseite
(Eigenwerte von $\Delta$) mit der geometrischen Seite (Längen geschlossener
Geodäten auf $\Gamma \backslash \HH$) verbindet. Sie ist das direkte Analogon
der Riemannschen expliziten Formel $\psi(x) = x - \sum_\rho x^\rho/\rho - \ldots$

\subsection{Das Weylsche Gesetz}

\begin{theorem}[Weylsches Gesetz für automorphe Spektren — bewiesen]
\label{satz:weyl}
Die Anzahl der Laplace-Eigenwerte $\lambda_j \leq T$ erfüllt
\[
N(T) = \#\{j : \lambda_j \leq T\} \sim \frac{\vol(\Gamma \backslash \HH)}{4\pi}\, T
\quad \text{für } T \to \infty.
\]
\end{theorem}

Das Weylsche Gesetz beweist die Existenz unendlich vieler Maass-Spitzenformen
für jede Kongruenzuntergruppe $\Gamma$.

\subsection{Die Selberg-\texorpdfstring{$\frac{1}{4}$}{1/4}-Vermutung}

\begin{conjecture}[Selberg, 1965 — offen]
\label{verm:selberg}
Für jede Kongruenzuntergruppe $\Gamma \leq \SL(2,\ZZ)$ und jeden Eigenwert
$\lambda > 0$ von $\Delta$ auf $L^2(\Gamma \backslash \HH)$ gilt
\[
\lambda \geq \frac{1}{4}.
\]
Äquivalent: alle Laplace-Eigenwerte entsprechen temperierten Darstellungen
($r_j \in \RR$, nicht rein imaginär).
\end{conjecture}

\begin{bemerkung}
Die bisher beste bekannte Schranke zur Vermutung~\ref{verm:selberg} ist
$\lambda_1 \geq \frac{1}{4} - \bigl(\frac{7}{64}\bigr)^2 \approx 0{,}2380$
von Kim und Sarnak (2003). Für $\Gamma = \SL(2,\ZZ)$ selbst legen numerische
Berechnungen $\lambda_1 \approx 91{,}14$ nahe (Booker-Strömbergsson), was mit
der Vermutung übereinstimmt.
\end{bemerkung}

% =============================================================================
\section{Das Langlands-Programm: Eine Vorschau}
\label{sec:langlands}
% =============================================================================

\subsection{Die Langlands-Korrespondenz}

Das Langlands-Programm, das Langlands 1967 in seinem Brief an Weil initiierte,
schlägt eine weitreichende Verallgemeinerung der Klassenkörpertheorie vor und
verknüpft:
\begin{itemize}
\item \textbf{Automorphe Darstellungen} $\pi$ von $G(\AA)$ für eine reduktive Gruppe $G/\QQ$;
\item \textbf{Galois-Darstellungen} $\rho: \Gal(\bar\QQ/\QQ) \to \hat{G}(\CC)$,
  wobei $\hat{G}$ die duale Langlands-Gruppe ist.
\end{itemize}
Für $G = \GL(1)$: dies ist Klassenkörpertheorie (Kronecker-Weber, Artin). Für $G = \GL(2)$:
dies umfasst die Shimura-Taniyama-Vermutung (bewiesen von Wiles-Taylor-Diamond-Conrad-Breuil).

\subsection{Automorphe $L$-Funktionen}

Zu jeder automorphen Darstellung $\pi = \otimes'_v \pi_v$ von $\GL(n,\AA)$ und
jeder endlichdimensionalen Darstellung $r: {}^L G \to \GL(N,\CC)$ der $L$-Gruppe
gehört eine $L$-Funktion
\[
L(s, \pi, r) = \prod_p L_p(s, \pi_p, r_p)
\]
als Euler-Produkt über Primzahlen. Für die Standarddarstellung $r = \mathrm{std}$
von ${}^L\GL(n) = \GL(n,\CC)$:
\[
L(s, \pi) = \prod_p \det\!\left(I - A_p\, p^{-s}\right)^{-1},
\]
wobei $A_p \in \GL(n,\CC)$ der \emph{Satake-Parameter} der unverzweigten Komponente
$\pi_p$ ist.

\subsection{Funktorialität}

\begin{conjecture}[Langlands-Funktorialität — teilweise bewiesen]
\label{verm:funktorialitaet}
Seien $G, H$ reduktive Gruppen über $\QQ$ und $\rho: {}^L G \to {}^L H$ ein
$L$-Homomorphismus. Dann existiert zu jeder automorphen Darstellung $\pi$ von
$G(\AA)$ eine automorphe Darstellung $\Pi$ von $H(\AA)$ — der \emph{funktori\-elle
Lift} — mit
\[
L(s, \Pi, r) = L(s, \pi, r \circ \rho)
\]
für alle Darstellungen $r$ von ${}^L H$.
\end{conjecture}

Diese Vermutung ist in mehreren wichtigen Fällen bewiesen: Basiswechsel für $\GL(n)$
(Arthur-Clozel), symmetrischer Quadratlift $\GL(2) \to \GL(3)$ (Gelbart-Jacquet),
und der symmetrische vierte Potenz-Lift $\GL(2) \to \GL(5)$ (Kim, 2003).

% =============================================================================
\section{Eisenstein-Reihen}
\label{sec:eisenstein}
% =============================================================================

\subsection{Klassische Eisenstein-Reihen}

Die \emph{Eisenstein-Reihe} vom Gewicht $k$ (gerade, $k \geq 4$) für $\SL(2,\ZZ)$
ist
\[
E_k(z) = 1 - \frac{2k}{B_k} \sum_{n=1}^\infty \sigma_{k-1}(n)\, q^n,
\quad q = e^{2\pi iz},
\]
wobei $B_k$ die $k$-te Bernoulli-Zahl ist und $\sigma_{k-1}(n) = \sum_{d|n} d^{k-1}$.
Beispielsweise gilt $E_4(z) = 1 + 240 \sum \sigma_3(n) q^n$ und
$E_6(z) = 1 - 504 \sum \sigma_5(n) q^n$.

\subsection{Nicht-holomorphe Eisenstein-Reihen}

Die \emph{nicht-holomorphe Eisenstein-Reihe} (Maass-Eisenstein-Reihe) für
$\Gamma = \SL(2,\ZZ)$ ist definiert durch
\[
E(z, s) = \sum_{\gamma \in \Gamma_\infty \backslash \Gamma} \im(\gamma z)^s
= \frac{1}{2} \sum_{\substack{(c,d)=1 \\ c,d \in \ZZ}} \frac{y^s}{|cz+d|^{2s}},
\quad \re(s) > 1,
\]
wobei $\Gamma_\infty = \bigl\{ \bigl(\begin{smallmatrix}1&n\\0&1\end{smallmatrix}\bigr) : n \in \ZZ \bigr\}$
der Stabilisator der Spitze $\infty$ ist.

\begin{theorem}[Meromorphe Fortsetzung der Eisenstein-Reihen — bewiesen]
\label{satz:eisenstein_forts}
Die Funktion $s \mapsto E(z,s)$ setzt sich zu einer meromorphen Funktion auf ganz $\CC$
fort. Sie besitzt einen einfachen Pol bei $s = 1$ mit Residuum $\vol(\Gamma \backslash \HH)^{-1}$
und ist holomorph auf $\re(s) = \frac{1}{2}$. Ferner erfüllt $E(z,s)$ die Funktionalgleichung
\[
E(z, s) = \varphi(s)\, E\!\left(z, 1-s\right),
\]
wobei
\[
\varphi(s) = \frac{\sqrt\pi\, \Gamma(s - \tfrac{1}{2})\,\zeta(2s-1)}{\Gamma(s)\,\zeta(2s)}
\]
die \emph{Streumatrix} ist (hier skalar, da $\Gamma$ eine Spitze hat).
\end{theorem}

\subsection{Fourier-Entwicklung von $E(z,s)$}

Die Fourier-Entwicklung von $E(z,s)$ lautet:
\begin{align*}
E(z,s) &= y^s + \varphi(s)\, y^{1-s} \\
&+ \frac{2\sqrt\pi\, y^{1/2}}{\Gamma(s)\,\zeta(2s)} \sum_{n \neq 0}
|n|^{s-1/2}\, \sigma_{1-2s}(|n|)\, K_{s-1/2}(2\pi|n|y)\, e^{2\pi i n x},
\end{align*}
wobei $\sigma_\nu(n) = \sum_{d|n} d^\nu$ und $K_\nu$ die Macdonald-Funktion ist.
Die beiden Terme $y^s + \varphi(s) y^{1-s}$ bilden den \emph{konstanten Term},
der zum stetigen Spektrum beiträgt.

% =============================================================================
\section{Kuspidales versus Eisenstein-Spektrum}
\label{sec:kuspidal_eis}
% =============================================================================

\subsection{Spektralzerlegung von \texorpdfstring{$L^2(\Gamma \backslash \HH)$}{L2}}

\begin{theorem}[Langlands-Spektralzerlegung — bewiesen, Langlands 1976]
\label{satz:langlands_zerlegung}
Für jede Kongruenzuntergruppe $\Gamma \leq \SL(2,\ZZ)$ gilt die orthogonale
direkte Summenzerlegung
\[
L^2(\Gamma \backslash \HH) = L^2_{\mathrm{cusp}}(\Gamma \backslash \HH) \oplus \CC \oplus L^2_{\mathrm{kont}}(\Gamma \backslash \HH),
\]
wobei:
\begin{enumerate}[label=(\roman*)]
\item $L^2_{\mathrm{cusp}}$ von den \emph{Maass-Spitzenformen} aufgespannt wird
  (mit diskreten Laplace-Eigenwerten $0 < \lambda_1 \leq \lambda_2 \leq \cdots$);
\item $\CC$ von der konstanten Funktion (Eigenwert $0$) aufgespannt wird;
\item $L^2_{\mathrm{kont}}$ das stetige Spektrum ist, das isometrisch durch die
  Eisenstein-Reihen $\bigl\{E(\cdot,\tfrac{1}{2}+it) : t \in \RR\bigr\}$ via
  die Eisenstein-Transformation
  \[
  \mathcal{E}: L^2(\RR_{>0}) \to L^2_{\mathrm{kont}},\quad
  \varphi \mapsto \frac{1}{4\pi}\int_{-\infty}^\infty \hat\varphi(t)\, E\!\left(\cdot, \tfrac{1}{2}+it\right) dt
  \]
  realisiert wird.
\end{enumerate}
\end{theorem}

\begin{bemerkung}
Langlands bewies diese Zerlegung in voller Allgemeinheit für beliebige reduktive
Gruppen über Zahlkörpern in seiner Monographie von 1976 \cite{langlands1976}.
Für $\GL(2)$ war sie bereits früher von Selberg (1956) und Roelcke (1956) aufgestellt
worden.
\end{bemerkung}

\subsection{Die Plancherel-Formel}

Die Zerlegung in Satz~\ref{satz:langlands_zerlegung} liefert die \emph{Plancherel-Formel}:
für $\phi, \psi \in L^2(\Gamma \backslash \HH)$ gilt
\[
\inner{\phi}{\psi} = \sum_j \hat\phi(r_j)\, \overline{\hat\psi(r_j)}
+ \frac{1}{4\pi} \int_{-\infty}^\infty \hat\phi(t)\, \overline{\hat\psi(t)}\, dt,
\]
wobei $\hat\phi(r_j) = \inner{\phi}{u_j}$ (Projektion auf Maass-Form $u_j$) und
$\hat\phi(t) = \inner{\phi}{E(\cdot, \frac{1}{2}+it)}$.

\subsection{Multiplizitäts-Eins-Satz}

\begin{theorem}[Multiplizitäts-Eins — bewiesen, Pjatezki-Shapiro 1979]
Jede irreduzible kusipidale automorphe Darstellung $\pi$ von $\GL(n,\AA)$
tritt im kuspidalen Spektrum $L^2_{\mathrm{cusp}}(\GL(n,\QQ) \backslash \GL(n,\AA))$
mit Multiplizität genau eins auf.
\end{theorem}

Dies ist ein fundamentales Strukturergebnis: Im Gegensatz zu kompakten Quotienten
oder allgemeinen reduktiven Gruppen ist das kusipidale Spektrum von $\GL(n)$ multiplizitätsfrei.

% =============================================================================
\section{Anwendungen}
\label{sec:anwendungen}
% =============================================================================

\subsection{Die Voronoi-Summenformel}

\begin{theorem}[Voronoi-Summation — bewiesen]
\label{satz:voronoi}
Sei $f : \RR_{>0} \to \CC$ eine glatte Funktion mit kompaktem Träger und
$\pi$ eine kusipidale automorphe Darstellung von $\GL(2,\AA)$ mit Fourier-Koeffizienten
$a(n)$ (für $n \geq 1$). Für $(a, q) = 1$ gilt:
\[
\sum_{n=1}^\infty a(n)\, e\!\left(\frac{an}{q}\right)\, f(n)
= \frac{1}{q} \sum_{\pm} \sum_{n=1}^\infty a(n)\, e\!\left(\frac{\pm \bar a n}{q}\right) \tilde f_\pm\!\!\left(\frac{n}{q^2}\right),
\]
wobei $\bar a \equiv a^{-1} \pmod{q}$ und $\tilde f_\pm$ explizite Integraltransformationen
von $f$ mit Besselfunktionen sind.
\end{theorem}

Die Voronoi-Summenformel ist ein mächtiges Werkzeug in der analytischen Zahlentheorie:
Sie transformiert Charaktersummen und Exponentialsummen in besser handhabbare duale
Summen. Sie hat Anwendungen bei Schranken für $L$-Funktionen und Charaktersummen.

\subsection{Subkonvexitätsschranken}

Ein zentrales Problem in der analytischen Theorie der $L$-Funktionen ist das Erhalten
von \emph{Subkonvexitätsschranken}: Schranken der Form
\[
L(\tfrac{1}{2} + it, \pi) \ll (|t| + 1)^{\frac{1}{4} - \delta}
\]
für ein $\delta > 0$. Die \emph{Konvexitätsschranke} (aus der Funktionalgleichung allein)
ergibt $\delta = 0$. Jede Verbesserung $\delta > 0$ heißt \emph{Subkonvexitätsschranke}.

\begin{theorem}[Weyl, Burgess, Michel-Venkatesh — verschiedene, bewiesen]
\begin{enumerate}[label=(\roman*)]
\item \textbf{Weyl (1921)}: Für die Riemann-Zetafunktion gilt
  $\zeta(\frac{1}{2}+it) \ll |t|^{1/6+\epsilon}$ (Subkonvexität im $t$-Aspekt).
\item \textbf{Burgess (1962)}: Für Dirichlet-$L$-Funktionen primitiver Charaktere
  $\chi \pmod{q}$ gilt $L(\frac{1}{2}, \chi) \ll q^{3/16+\epsilon}$ (Subkonvexität
  im $q$-Aspekt).
\item \textbf{Michel-Venkatesh (2010)}: Subkonvexität für $\GL(1)$- und $\GL(2)$-$L$-Funktionen
  in allen Aspekten gleichzeitig, mit Anwendung auf die Gleichverteilung der Heegner-Punkte.
\end{enumerate}
\end{theorem}

\subsection{Der Sato-Tate-Satz}

\begin{theorem}[Sato-Tate — bewiesen, Taylor u.\,a.\ 2006--2011]
\label{satz:sato_tate}
Sei $E/\QQ$ eine elliptische Kurve ohne komplexe Multiplikation, und schreibe
$a_p = p + 1 - \#E(\FF_p) = 2\sqrt{p}\cos(\theta_p)$ mit $\theta_p \in [0,\pi]$.
Dann sind die Winkel $\theta_p$ gleichverteilt bezüglich des Sato-Tate-Maßes
\[
d\mu_{\mathrm{ST}} = \frac{2}{\pi}\sin^2(\theta)\, d\theta.
\]
Äquivalent: für jedes $[a,b] \subset [0,\pi]$ gilt
\[
\frac{\#\{p \leq X : \theta_p \in [a,b]\}}{\#\{p \leq X\}} \longrightarrow \frac{2}{\pi} \int_a^b \sin^2\theta\, d\theta
\quad \text{für } X \to \infty.
\]
\end{theorem}

Der Beweis von Clozel, Harris, Shepherd-Barron und Taylor beruht auf:
(1) den symmetrischen Potenz-Lifts $\mathrm{Sym}^n(f_E)$ für die Hecke-Eigenform
$f_E$ zu $E$ (via Modularität), (2) potentieller Automorphie dieser symmetrischen
Potenzen mittels eines neuartigen Deformationsarguments, und (3) der analytischen
Fortsetzung und Nicht-Verschwindung der resultierenden $L$-Funktionen auf $\re(s) = 1$.

% =============================================================================
\section{Neuere Entwicklungen}
\label{sec:neueres}
% =============================================================================

\subsection{Langlands-Funktorialität für \texorpdfstring{$\GL(2) \to \GL(n)$}{GL(2) nach GL(n)}}

Die spektakulärsten neueren Fortschritte zur Funktorialität betreffen die
symmetrischen Potenz-Lifts von $\GL(2)$ nach $\GL(n)$:

\begin{theorem}[Symmetrische Potenz-Lifts — teilweise bewiesen]
\label{satz:sym_lift}
Sei $\pi$ eine kusipidale automorphe Darstellung von $\GL(2,\AA)$, die nicht vom
diedralen, tetraedrischen oder oktaedrischen Typ ist. Dann gilt:
\begin{enumerate}[label=(\roman*)]
\item $\mathrm{Sym}^2 \pi$ existiert als automorphe Darstellung von $\GL(3,\AA)$
  (Gelbart-Jacquet, 1978 — bewiesen);
\item $\mathrm{Sym}^3 \pi$ existiert als automorphe Darstellung von $\GL(4,\AA)$
  (Kim-Shahidi, 2002 — bewiesen);
\item $\mathrm{Sym}^4 \pi$ existiert als automorphe Darstellung von $\GL(5,\AA)$
  (Kim, 2003 — bewiesen);
\item $\mathrm{Sym}^n \pi$ für $n \geq 5$: offen im Allgemeinen (Vermutung).
\end{enumerate}
\end{theorem}

Die Existenz von $\mathrm{Sym}^4 \pi$ hat das wichtige Korollar, dass die
Ramanujan-Vermutung bis auf $|a_p| \leq 2p^{7/64}$ gilt (Kim-Sarnak), was die
triviale Schranke $p^{k/2}$ in Richtung der vermuteten $p^{(k-1)/2}$ verbessert.

\subsection{Das Fundamentallemma}

\begin{theorem}[Fundamentallemma — bewiesen, Ngô Bảo Châu, 2010, Fields-Medaille]
\label{satz:fundamentallemma}
Sei $G$ eine zusammenhängende reduktive Gruppe über einem nicht-archimedischen
lokalen Körper $F$, und $H$ eine endoskopische Gruppe von $G$. Seien $f_H \in \mathcal{H}(H(F))$
und $f \in \mathcal{H}(G(F))$ passende Funktionen. Dann gelten die von der Endoskopietheo\-rie
vorhergesagten orbitalen Integralidentitäten:
\[
\text{SO}(\gamma_H, f_H) = \Delta(\gamma_H, \gamma)\, \text{O}(\gamma, f)
\]
für passende stark regulär halbeinfache Elemente $\gamma_H \in H(F)$ und
$\gamma \in G(F)$, wobei $\Delta(\gamma_H,\gamma)$ der Langlands-Shelstad-Transferfaktor ist.
\end{theorem}

Das Fundamentallemma wurde von Langlands und Shelstad 1983 vermutet und in
verschiedenen Fällen bewiesen, bevor Ngô einen vollständigen Beweis mittels der
Geometrie der Hitchin-Faserungen über endlichen Körpern gab. Ngôs Beweis reduziert
die Identität auf das Zählen von Fixpunkten des Frobenius auf bestimmten algebraischen
Varietäten (Hitchin-Fasern) unter Verwendung des Zerlegungssatzes von
Beilinson-Bernstein-Deligne. Die im Jahr 2010 an Ngô Bảo Châu verliehene Fields-Medaille
würdigte dies als „den Beweis einer der wichtigsten Vermutungen des Langlands-Programms."

\subsection{Das geometrische Langlands-Programm}

Ein verwandtes, aber eigenständiges Programm, das von Drinfeld und Laumon (1987)
initiiert und von Frenkel, Ben-Zvi, Nadler u.\,a. weiterentwickelt wurde, formuliert
ein Analogon der Langlands-Korrespondenz für algebraische Kurven über $\CC$:
\begin{itemize}
\item Automorphe Darstellungen $\to$ $\mathcal{D}$-Moduln auf dem Modulstack der
  $G$-Bündel $\mathrm{Bun}_G$;
\item Galois-Darstellungen $\to$ flache $\hat{G}$-Zusammenhänge (oder, in der
  quantischen Version, $\hat{G}$-lokale Systeme).
\end{itemize}
Das geometrische Programm ist in vielerlei Hinsicht zugänglicher (man kann die
volle Kraft der algebraischen Geometrie, perverser Garben und $\infty$-Kategorien
nutzen) und hat Rückwirkungen auf das klassische Programm.

% =============================================================================
\section*{Schlussbetrachtung}
% =============================================================================

Automorphe Formen stehen an der Schnittstelle einiger der tiefsten Strömungen
der modernen Mathematik. Die Spektraltheorie von $\Delta$ auf arithmetischen
Quotienten, von Selberg (bewiesen, 1956) entwickelt, bietet ein unendlichdimensionales
Analogon der klassischen Fourier-Analysis, dessen Reichtum noch bei weitem nicht
erschöpft ist. Das Langlands-Programm — mit seinen bewiesenen Komponenten
(Ngôs Fundamentallemma, symmetrische Quadrat-, Kubik- und Vierter-Potenz-Lifts,
Sato-Tate, Voronoi-Summation) und seinen offenen Teilen (allgemeine Funktorialität,
die Selberg-$\tfrac{1}{4}$-Vermutung, die vollständige Ramanujan-Vermutung) —
bleibt eine der zentralen Forschungsfronten der Mathematik im einundzwanzigsten Jahrhundert.

Das Modul \texttt{automorphic\_forms.py} im Specialist Mathematics Project
implementiert die Kernstrukturen: \texttt{AutomorphicForm}, \texttt{SpectralDecomposition},
\texttt{WhittakerModel}, \texttt{GlobalLFunction}, \texttt{FunctorialLift},
\texttt{RamanujanConjecture} und \texttt{SatakeTransform}.

% =============================================================================
\begin{thebibliography}{99}

\bibitem{selberg1956}
A.~Selberg,
\textit{Harmonic analysis and discontinuous groups in weakly symmetric Riemannian
spaces with applications to Dirichlet series},
J.~Indian Math.\ Soc.\ \textbf{20} (1956), 47--87.

\bibitem{langlands1970}
R.~P.~Langlands,
\textit{Problems in the theory of automorphic forms},
Lectures in Modern Analysis and Applications III,
Lecture Notes in Mathematics \textbf{170}, Springer, 1970, S.~18--61.

\bibitem{langlands1976}
R.~P.~Langlands,
\textit{On the Functional Equations Satisfied by Eisenstein Series},
Lecture Notes in Mathematics \textbf{544}, Springer, 1976.

\bibitem{iwaniec_kowalski}
H.~Iwaniec und E.~Kowalski,
\textit{Analytic Number Theory},
American Mathematical Society Colloquium Publications \textbf{53}, AMS, 2004.

\bibitem{bump}
D.~Bump,
\textit{Automorphic Forms and Representations},
Cambridge Studies in Advanced Mathematics \textbf{55}, Cambridge University Press, 1997.

\bibitem{cogdell}
J.~W.~Cogdell,
\textit{Lectures on $L$-functions, Converse Theorems, and Functoriality for $\GL_n$},
in: \textit{Lectures on Automorphic $L$-Functions},
Fields Institute Monographs \textbf{20}, AMS, 2004, S.~1--96.

\bibitem{ngo2010}
B.~C.~Ngô,
\textit{Le lemme fondamental pour les algèbres de Lie},
Publications Mathématiques de l'IHÉS \textbf{111} (2010), 1--169.

\bibitem{kim2003}
H.~H.~Kim,
\textit{Functoriality for the exterior square of $\GL_4$ and the symmetric fourth of $\GL_2$},
J.~Amer.\ Math.\ Soc.\ \textbf{16} (2003), 139--183.

\bibitem{taylor2008}
R.~Taylor,
\textit{Automorphy for some $\ell$-adic lifts of automorphic mod $\ell$ Galois representations, II},
Publications Mathématiques de l'IHÉS \textbf{108} (2008), 183--239.

\bibitem{michel_venkatesh}
P.~Michel und A.~Venkatesh,
\textit{The subconvexity problem for $\GL_2$},
Publications Mathématiques de l'IHÉS \textbf{111} (2010), 171--271.

\end{thebibliography}

\end{document}
